\chapter{Alcances del sistema}
\label{cap:alcances}

El proyecto comprenderá el análisis del proceso de venta y reserva de
paquetes turísticos de Vizcacha Tours y el desarrollo de un sistema web que
permita administrar paquetes, programar salidas, consultar la disponibilidad
de cupos, registrar clientes y pasajeros, elaborar cotizaciones, gestionar
reservas, registrar anticipos y saldos, administrar proveedores y generar
reportes. La prioridad será verificar los cupos y registrar correctamente las
reservas; las demás funciones se incorporarán como apoyo a ese proceso
principal. El sistema será utilizado por el personal de la agencia y por los
clientes, que accederán con una cuenta propia para consultar la información
de los paquetes y solicitar sus reservas.

\section{Módulos propuestos}

\begin{itemize}
  \item \textbf{Paquetes turísticos:} registro de destinos, descripción,
        duración, itinerario, precios y servicios incluidos.
  \item \textbf{Salidas y disponibilidad:} programación de fechas, capacidad
        de pasajeros y cálculo de cupos disponibles.
  \item \textbf{Clientes y pasajeros:} registro de datos necesarios para
        atender y dar seguimiento a las solicitudes.
  \item \textbf{Cotizaciones:} preparación y consulta de propuestas según el
        paquete solicitado y sus condiciones.
  \item \textbf{Reservas:} creación, consulta, confirmación, modificación y
        cancelación con validación de cupos. La reserva puede originarse en la
        atención del personal o en la solicitud del cliente; en ambos casos se
        registra en estado pendiente y requiere la confirmación del personal.
  \item \textbf{Pagos:} registro de anticipos y saldos, con pago en efectivo y
        simulación de banca móvil o código QR, sin integración con pasarelas
        bancarias reales. El anticipo no puede ser menor al 30\% del total de
        la reserva, la totalidad también puede pagarse por la web y el pago
        registrado por el cliente queda pendiente de validación por parte del
        personal. La devolución del dinero se rige por tramos según los días
        de anticipación de la cancelación y su ejecución queda fuera del
        sistema.
  \item \textbf{Acceso del cliente:} registro de cuenta, con el correo y el
        documento obligatorios y únicos, consulta de paquetes,
        salidas y cupos disponibles, solicitud de reserva, seguimiento de sus
        propias reservas, registro de pagos y pagos simulados, con acceso
        restringido a la información que le corresponde.
  \item \textbf{Proveedores:} datos básicos de transporte, hospedaje y guías
        vinculados a los paquetes.
  \item \textbf{Reportes y usuarios:} consultas por fecha, ocupación, reservas
        y accesos según el rol de cada usuario, sea trabajador o cliente.
\end{itemize}

Para facilitar su uso, se procurará una interfaz sencilla y adaptable a
distintos tamaños de pantalla.

\section{Canales de atención y confirmación}

El sistema contempla dos canales de atención que desembocan en el mismo
proceso de reserva. El primero es la atención presencial y telefónica, en la
que el personal registra la consulta, prepara la cotización y confirma la
reserva. El segundo es el acceso directo del cliente al sistema web, en el que
consulta los paquetes, las salidas y los cupos disponibles, y remite su
solicitud de reserva junto con el registro de su pago.

En ambos casos la reserva se crea en estado pendiente y requiere la validación
y la confirmación del personal, con la misma verificación de disponibilidad y
el mismo registro de auditoría. La confirmación exige el respaldo del anticipo
validado, y las solicitudes del cliente que no son confirmadas dentro del
plazo de 24 horas quedan vencidas y liberan automáticamente los cupos
reservados. De esta manera, el acceso directo del
cliente amplía la cobertura del proceso sin duplicar reglas ni alterar las
tareas de confirmación, seguimiento y control que corresponden al personal.

\section{Avisos al cliente}

Los cambios en la fecha o en la capacidad de una salida Programada y las
confirmaciones de reserva se comunicarán al cliente afectado a través de las
pantallas del sistema, lo que reduce la demora en comunicar estas novedades y
mantiene informada a la persona que viaja.
